\newcommand{\objone}{
\begin{enumerate}
    \item Accept a range of numbers (start and end) as command line arguments to the shell script.
    \item Identify and check each number within the specified range to determine if it is a prime number.
    \item Print all prime numbers found within the given range to the terminal.
\end{enumerate}
}

\newcommand{\objtwo}{\\
The objective of this assignment is to write a shell script that generates and displays Pascal’s Triangle up to a specified number of rows. This script demonstrates the use of loops, conditional statements, and arithmetic operations in shell programming. Pascal’s Triangle is a mathematical structure used to represent binomial coefficients, and this exercise helps in understanding both combinatorics and basic shell scripting concepts.
}

\newcommand{\objthree}{\\
To practice file manipulation and command-line arguments in shell scripting by writing a script that removes all lines containing a specific keyword from a given file.
This exercise helps in understanding:
\begin{itemize}
    \item How to pass and use command-line arguments in shell scripts.
    \item How to use $grep$, $sed$, or $awk$ for pattern matching and text processing.
    \item Basic file handling operations using shell scripting.
\end{itemize}
}

\newcommand{\objfour}{\\
To develop a shell script that extracts and displays a specific range of lines from a text file using command-line arguments.
This exercise helps in understanding:
\begin{itemize}
    \item Command-line argument handling in shell scripting.
    \item Usage of line processing tools like $sed$.
    \item File reading and conditional execution in bash scripts.
\end{itemize}
}

\newcommand{\objfive}{\\
To create a shell script that filters and displays only those lines from a file that begin with a vowel. The filename is passed as a command-line argument.
This enhances understanding of:
\begin{itemize}
    \item Command-line argument usage.
    \item Regular expressions in text processing.
    \item File filtering using $grep$ in shell scripting.
\end{itemize}
}

\newcommand{\objsix}{\\
To write a shell script that processes multiple command-line arguments, identifies whether each argument is a file or directory, and provides detailed information including the number of lines if it is a file.
This task helps in learning:
\begin{itemize}
    \item Handling multiple command-line arguments in shell scripts.
    \item Using conditional statements to distinguish between files and directories.
    \item Counting lines in a file using command-line utilities.
\end{itemize}
}

\newcommand{\objseven}{\\
To develop a shell script that accepts a folder name as a command-line argument and identifies all empty subfolders within it, saving their names to an output file.
This exercise aims to strengthen understanding of:
\begin{itemize}
    \item Directory traversal using shell scripting.
    \item Detecting empty directories.
    \item Command-line argument handling and file output operations.
\end{itemize}
}

\newcommand{\objeight}{\\
The objective of this assignment is to write a shell script that simulates the functionality 
of the 'wc' (word count) command, which displays the number of lines, words, and characters 
in a given text file.

This assignment aims to:
\begin{itemize}
\item Develop the ability to work with file input and text processing in shell scripting.
\item Reinforce the use of commands like 'wc', 'cat', 'tr', and 'wc -l', 'wc -w', 'wc -c'.
\item Practice the implementation of command substitution and argument validation.
\item Strengthen scripting logic to create alternatives to standard Linux utilities.
\end{itemize}
}

\newcommand{\objnine}{\\
 To write a shell script that implements the Bubble Sort algorithm on an array of numbers. The sorting order is determined by a command-line argument: if \texttt{1} is provided, the array is sorted in ascending order; if \texttt{2} is provided, the array is sorted in descending order.
This assignment helps to understand:
\begin{itemize}
    \item Implementation of sorting algorithms in shell scripts.
    \item Handling arrays and loops in shell scripting.
    \item Command-line argument parsing and conditional logic.
\end{itemize}
}

\newcommand{\objten}{\\
The objective of this assignment is to write a shell script that prints the Fibonacci series up to a given number, where the range (limit) is provided as a command line argument. The script should validate the input and generate the series correctly up to the specified value.
This assignment is designed to:
\begin{enumerate}
    \item Strengthen skills in shell scripting using command line arguments.

    \item Develop logic-building ability by implementing iterative control structures (like while or for loops).

    \item Enhance understanding of variable handling, arithmetic operations, and input validation in shell scripts.

    \item Practice creating user-interactive scripts that accept dynamic inputs.
\end{enumerate}
}

\newcommand{\objeleven}{\\
To write a shell script that finds the second highest number from a list of user-provided numbers without using any sorting technique.
This assignment helps in understanding:
\begin{itemize}
    \item Array and variable manipulation in shell scripting.
    \item Use of conditional logic to track maximum and second maximum values.
    \item Basic algorithm design without relying on built-in sorting utilities.
\end{itemize}
}

\newcommand{\objtwelve}{\\
To write a shell script that performs binary search on a set of sorted data elements to find the position of a specified target value.
This assignment aims to enhance understanding of:
\begin{itemize}
    \item Implementing search algorithms in shell scripts.
    \item Efficiently accessing and manipulating array elements.
    \item Using control structures like loops and conditional statements for algorithmic logic.
\end{itemize}
}

\newcommand{\objthirteen}{\\
The objective of this assignment is to write a shell script that analyzes a given text file and prepares a word frequency report, showing how many times each unique word appears in the file.
}

\newcommand{\objfourteen}{\\
To write a shell script that accepts a text file as input and removes all duplicate (identical) lines, retaining only the unique lines in the output.

This assignment helps in understanding:
\begin{itemize}
    \item File handling and reading in shell scripts.
    \item Use of text processing tools like \texttt{sort} and \texttt{uniq}.
    \item Practical applications of command-line utilities for data cleaning.
\end{itemize}
}

\newcommand{\objfifteen}{\\
The objective of this assignment is to write a shell script that reads data from a source file and creates a new file by copying the content while excluding all vowels (both uppercase and lowercase).
}

\newcommand{\objsixteen}{\\
To write a shell script that calculates the roots of a quadratic equation of the form $ax^2 + bx + c = 0$, where $a \ne 0$. The script should handle all possible cases: real and distinct roots, real and equal roots, and complex (imaginary) roots.

This assignment aims to develop:
\begin{itemize}
    \item Mathematical computation using shell scripting.
    \item Conditional branching to handle different types of roots based on the discriminant.
    \item Command-line input handling and basic floating-point arithmetic.
\end{itemize}
}

\newcommand{\objseventeen}{\\
To write a shell script that performs Selection Sort on an array of numbers to arrange them in descending order.

This assignment helps in understanding:
\begin{itemize}
    \item Implementation of sorting algorithms using shell scripting.
    \item Array manipulation and indexing in bash.
    \item Use of loops and conditional statements to perform comparisons and swaps.
\end{itemize}

}

\newcommand{\objeighteen}{\\
The objective of this assignment is to write a shell script that calculates the Greatest Common Divisor (GCD) of two integers provided by the user (via input or command line arguments).
}

\newcommand{\objnineteen}{\\
 To write a shell script that implements a basic calculator using a menu-driven approach. The script should perform addition, subtraction, multiplication, and division based on the user's choice, which is provided as a command-line argument.

This assignment helps in understanding:
\begin{itemize}
    \item Command-line argument parsing in shell scripts.
    \item Implementation of arithmetic operations using shell scripting.
    \item Use of conditional statements and control structures to build menu-driven programs.
\end{itemize}
}

\newcommand{\objtwenty}{\\
The objective of this assignment is to write a shell script that checks whether a given number and its reverse are the same (i.e., whether the number is a palindrome), where the number is passed as a command line argument.
This assignment is intended to:
\begin{itemize}
    \item Strengthen the understanding of command line argument handling in shell scripts.
    \item Practice control flow statements and loops for number manipulation.
    \item Enhance logic-building skills by implementing digit extraction and reversal.
    \item Develop familiarity with conditional checking and basic input validation.
    \item Encourage writing user-friendly and functional shell scripts for numeric operations.
\end{itemize}
}

\newcommand{\objtwentyone}{\\
To write a shell script that converts the content of a file from lowercase to uppercase or from uppercase to lowercase using a menu-driven approach.

This assignment helps in understanding:
\begin{itemize}
    \item File reading and writing in shell scripting.
    \item Text transformation using commands like \texttt{tr}.
    \item Implementation of menu-driven logic using control structures.
    \item Command-line interaction and user input handling.
\end{itemize}
}

\newcommand{\objtwentytwo}{\\
To write a shell script that implements the Bubble Sort algorithm to sort an array of numbers.
This assignment helps in understanding:
\begin{itemize}
    \item The logic and working of the Bubble Sort algorithm.
    \item Array handling and manipulation in shell scripting.
    \item Use of loops and conditional statements to perform comparisons and swaps.
\end{itemize}
}

\newcommand{\objtwentythree}{\\
To write a shell script that prints a specific star pattern using nested loops.
This assignment helps in understanding:
\begin{itemize}
    \item Looping constructs (for/while) in shell scripting.
    \item Printing formatted output to the terminal.
    \item Logical thinking for constructing symmetrical patterns using loops.
\end{itemize}
}

\newcommand{\objtwentyfour}{\\
To write a shell script that takes a number as a command-line argument and prints all of its prime factors.

This assignment helps in understanding:
\begin{itemize}
    \item Command-line argument handling in shell scripting.
    \item Implementation of basic mathematical algorithms (prime factorization).
    \item Use of loops and conditional statements for number theory operations.
\end{itemize}
}

\newcommand{\objtwentyfive}{\\
To write a shell script that takes a word as input from the user and calculates the frequency of that word in a specified file.
This assignment helps in understanding:
\begin{itemize}
    \item User input handling in shell scripts.
    \item File reading and text processing.
    \item Usage of command-line tools like \texttt{grep}, \texttt{wc}, and pattern matching to count word occurrences.
\end{itemize}

}

\newcommand{\objtwentysix}{\\
The objective of this assignment is to write a shell script that checks whether a given word 
is a palindrome, where the word is provided as a command line argument, and its length 
must be at least 4 characters.
\begin{itemize}
This assignment aims to:
- Reinforce handling of command line arguments and string operations in shell scripting.
- Practice input validation based on word length.
- Strengthen logic implementation for reversing and comparing strings.
- Develop scripts that perform basic text analysis tasks effectively.
\end{itemize}
}

\newcommand{\objtwentyseven}{\\
To write a shell script that takes a filename as input at runtime and counts the total number of words, lines, and characters in the file.

This assignment helps in understanding:
\begin{itemize}
    \item Command-line input handling in shell scripting.
    \item File processing using standard Unix commands.
    \item Utilization of tools like \texttt{wc} to analyze text files.
\end{itemize}

}

\newcommand{\objtwentyeight}{\\
The objective of this assignment is to write a shell script that executes automatically 
at the time of user login and displays an appropriate greeting such as 
"Good Morning", "Good Afternoon", "Good Evening", or "Good Night" 
based on the current system time.

This assignment aims to:
\begin{itemize}
\item Demonstrate the use of system time and conditional logic in shell scripting.
\item Familiarize with customizing user login experience via shell configuration files.
\item Enhance understanding of time-based automation using shell scripts.
\item Develop practical skills in managing user environment behavior during login.
\end{itemize}
}

\newcommand{\objtwentynine}{\\
To write a shell script that implements the Sequential Search (Linear Search) algorithm to find a number in a list of numbers.
This assignment helps in understanding:
\begin{itemize}
    \item Implementation of searching algorithms in shell scripting.
    \item Array handling and user input processing.
    \item Use of loops and conditional statements to search for a specific value.
\end{itemize}
}

\newcommand{\objthirty}{\\
The objective of this assignment is to write a shell script that implements the binary 
search algorithm to find a given number in a sorted list of numbers provided by the user.

This assignment aims to:
\begin{itemize}
\item Strengthen understanding of fundamental searching algorithms using shell scripting.
\item Practice handling arrays, user input, and function implementation in Bash.
\item Demonstrate the use of loop structures and conditional statements for logic control.
\item Enhance problem-solving skills using algorithmic thinking in shell scripts.
\end{itemize}
}

\newcommand{\objthirtyone}{\\
The objective of this assignment is to write a shell script to calculate the factorial 
of a given non-negative integer supplied through the command line or user input.

This assignment aims to:
\begin{enumerate}
\item Practice the use of loops and arithmetic operations in shell scripting.
\item Reinforce the concept of iterative problem-solving using control structures.
\item Demonstrate input validation and error handling in scripts.
\item Enhance logic development skills through mathematical computation in shell.
\end{enumerate}
}

\newcommand{\objthirtytwo}{\\
To write a shell script that reverses a given string and checks whether the reversed string is the same as the original, thereby determining if it is a palindrome.

This assignment helps in understanding:
\begin{itemize}
    \item String manipulation techniques in shell scripting.
    \item Use of command-line tools like \texttt{rev} or manual reversal logic.
    \item Implementation of conditional checks to compare strings.
\end{itemize}
}

\newcommand{\objthirtythree}{\\
The objective of this assignment is to write a menu-driven shell script to manage 
employee records by allowing insertion of employee details (ID, name, designation, salary) 
into a data file, displaying all records, and retrieving salary information for a specific employee.
}